\section{Method-specific tuning and an independent replication}\label{r11:robustness}
The preceding holdout answers a precisely frozen comparison, but its scalar learning-rate search leaves an important design question: does the ordering survive a more consequential method-specific tuning effort? We answer this question with a new experiment, rather than choosing new hyperparameters on the already observed holdout or suppressing the earlier results. The objective, diffusion, action box, horizon, equal actor/critic initializations, author solver commit, and 128-date deployed policy are unchanged. The new protocol fixes tuning seeds 1600 and 1601 and holdout seeds 1820--1831 before execution. None of these is an R9 holdout seed.

\subsection{Prospectively specified equal tuning opportunities}
For each dimension and method, twelve initial configurations and two prospectively specified extension configurations are each trained for five nominal seconds on two tuning seeds. Thus each method receives 140 nominal training seconds per dimension. Completed-update overshoot, initialization, validation, and holdout clocks are recorded separately. The methods receive equal nominal training opportunity, not equal numbers of gradient operations or an assertion that their entire hyperparameter spaces have been optimized.

The shared learning-rate multiplier set expands from $\{1/3,1,3\}$ to $\{1/3,1,3,9,27\}$. NBO additionally varies evaluation updates per actor improvement and transition particles. SOC-MartNet additionally varies the adversary learning rate, the number of sinusoidal test outputs, the solver's $J/K$ parameters, and the initial and updating multiplier parameters. The complete twelve-configuration lists appear in the frozen executable protocol; unchanged settings and unsuccessful candidates are not omitted. The pinned author training implementation is unmodified. Method order is randomized within tuning seed and candidate, and actor and critic initial hashes must agree within each holdout pair.

After the first twelve configurations, selection minimizes the mean validation cost over the two tuning seeds, with an index-based tie rule. A selected multiplier at the current high or low boundary is extended by a factor of three in that direction. If the multiplier is interior, the additional slot varies transition-particle count for NBO or adversary rate for SOC. Selection is repeated before the second extension. These rules are specified before execution and do not inspect the new holdout. Failed tuning trials receive infinite selection score while their raw records are preserved. Validation uses 1,024 common Brownian paths per tuning seed, and holdout evaluation uses 4,096 new common paths per seed. All recorded outcomes in this execution are finite, and all 168 tuning trials and 36 holdout pairs are retained.

\begin{table}[t]\centering\small
\caption{Selected configurations and actual tuning-training opportunity}\label{r11:tab:tuning}
\begin{tabular}{rllrr}
\toprule
$d$ & NBO selected & SOC selected & NBO sec. & SOC sec. \\
\midrule
8 & 3 ($m=9$) & 11 ($m=3$) & 140.15 & 140.29 \\
16 & 12 ($m=3$) & 2 ($m=3$) & 140.09 & 140.18 \\
32 & 1 ($m=1$) & 3 ($m=9$) & 140.17 & 140.31 \\
\bottomrule
\end{tabular}

\par\smallskip\footnotesize Configuration indices are zero based and bind to the complete frozen lists and extension records. $m$ is the actor/critic learning-rate multiplier. Actual seconds sum all 28 tuning fits per method and dimension, including completed-update overshoot but excluding separately retained validation and setup.
\end{table}

All six selected learning-rate multipliers are interior to the enlarged tested multiplier set. At $d=8$, NBO selects $m=9$, while SOC selects $m=3$ with initial multiplier and multiplier step both 30. At $d=16$, NBO selects six evaluation steps and sixteen transition particles with $m=3$; SOC selects its default method-specific parameters with $m=3$. At $d=32$, NBO selects $m=1$ and SOC selects $m=9$, with the respective default remaining parameters. The large learning-rate boundary identified in the earlier report is therefore tested beyond its old limit. Interior selection on this finite set does not establish global tuning optimality or exclude improvements along an unsearched interaction. The experiment tests a finite, disclosed method-specific tuning budget.

\subsection{Independent holdout and protocol sensitivity}
The selected configurations are fixed before their twelve new holdout pairs are run at a ten-second nominal training budget. The primary test is again the one-sided exact sign test, now with Bonferroni correction across the \emph{three new} predeclared panels. This does not change the six-panel correction of the earlier experiment. Because the historical evidence motivated the new experiment, the new study is a prospective post-review replication, not a retroactive preregistration of the original benchmark design.

\begin{table}[t]\centering\scriptsize
\caption{New holdout after equal-budget method-specific tuning}\label{r11:tab:robustness}
\begin{tabular}{rrrrrrrr}
\toprule
$d$ & NBO & SOC & LQ & Mean diff. & Median diff. & Wins & Adj. $p$ \\
\midrule
8 & 0.5542 & 0.7928 & 0.5125 & -0.2386 & -0.1967 & 11/12 & 0.00952 \\
16 & 0.4547 & 0.8067 & 0.1741 & -0.3520 & -0.1212 & 8/12 & 0.58154 \\
32 & 0.3824 & 0.4586 & 0.0852 & -0.0762 & -0.0427 & 7/12 & 1.00000 \\
\bottomrule
\end{tabular}

\par\smallskip\footnotesize Smaller realized cost is better. Differences are NBO minus SOC, using paired audit paths and twelve independently seeded training-and-audit pipelines. The last column applies the prospectively fixed three-panel sign-test correction. Costs are sampled policy evaluations, not outward deterministic certificates. The LQ column retains the independently specified clipped feedback under the unchanged nonconvex objective.
\end{table}

The NBO mean differences are $-.23863$, $-.35199$, and $-.07623$ in dimensions 8, 16, and 32. The corresponding paired win counts are 11, 8, and 7 out of twelve, with adjusted sign-test values $.00952$, $.58154$, and 1. Only the eight-dimensional panel establishes the declared robust seed-level ordering at level $.05$. In dimension 16, four seeds again favor SOC. In dimension 32, the expanded tuning yields a substantially smaller mean difference and no significant directional ordering, although the earlier frozen experiment favored NBO in all twelve seeds. The improved SOC result is consistent with its selected $m=9$, which the old scalar grid excluded; this association is not an isolated causal estimate of that hyperparameter because tuning, validation seeds, and holdout seeds all differ.

The independently specified clipped-feedback policy has lower mean realized cost than either trained pipeline in each new panel, as it did in the previous experiment. Its favorable performance does not make its cost a lower bound for a minimization problem, and it does not prove global optimality. These outcomes distinguish three questions: whether a configured NBO pipeline improves on a configured external solver, whether the ordering is robust to a changed tuning protocol, and whether either neural policy has small absolute error. The new experiment supplies direct evidence for the second question and retains all evidence relevant to the first; Section~\ref{r11:accuracy} supplies a separate exact-reference experiment for the third. Favorable absolute-accuracy results in another model are not used to erase the sixteen- or thirty-two-dimensional sign-test outcomes here.

The evidence pipeline collects with an unconditional completion policy. Every planned cell has explicit status before dependency installation; successful outputs, failed outputs, and missing-cell metadata are versioned before aggregate validity is checked. Separate immutable commits identify the frozen study source and completed raw results. No branch name alone is used as a review target, and the 36 new pairs neither replace nor pool selectively with the 72 earlier pairs.
